\documentclass[12pt,a4paper]{article}

\usepackage{fontspec}
\usepackage{polyglossia}
\setdefaultlanguage{turkish}
\setmainfont{Times New Roman}
\usepackage[a4paper,margin=2.5cm]{geometry}
\usepackage{setspace}
\usepackage{parskip}
\usepackage{titlesec}
\usepackage{hyperref}
\usepackage{enumitem}
\usepackage{amsmath,amssymb}
\usepackage{booktabs}
\usepackage{tabularx}
\usepackage{longtable}
\usepackage{array}
\usepackage{pgfgantt}
\usepackage{xcolor}
\usepackage{fancyhdr}

\hypersetup{colorlinks=true,linkcolor=black,urlcolor=blue,citecolor=blue,pdftitle={Universite Tarzi Proposal TR}}
\setlength{\headheight}{15pt}
\pagestyle{fancy}
\fancyhf{}
\lhead{Research Proposal}
\rhead{Aging-Aware Adaptive Scheduling}
\cfoot{\thepage}

\definecolor{deepblue}{RGB}{0,51,102}
\definecolor{accentblue}{RGB}{0,102,204}
\titleformat{\section}{\large\bfseries\color{deepblue}}{\thesection.}{0.6em}{}
\titleformat{\subsection}{\normalsize\bfseries\color{accentblue}}{\thesubsection.}{0.6em}{}

\begin{document}

\begin{titlepage}
\begin{center}
{\Large \textbf{Izmir Yuksek Teknoloji Enstitusu}}\\[0.8em]
{\large Muhendislik Fakultesi}\\
{\large Bilgisayar Muhendisligi}\\[2.5em]

{\LARGE \textbf{TEZ / PROJE ONERISI}}\\[1.0em]
{\Large \textbf{Yapay Zeka Destekli Aging-Aware Adaptive Time Quantum Scheduling}}\\[0.4em]
{\large Cross-Layer Reliability Feedback ile Guvenilir ve Verimli Scheduler Tasarimi}\\[2.0em]

\begin{tabular}{>{\bfseries}p{4.8cm}p{8.8cm}}
Ogrenci Adi Soyadi: & [Ad Soyad] \\
Ogrenci Adi Soyadi: & Omer Faruk Teker \\
Ogrenci Numarasi: & Belirtilecek \\
Danisman: & Belirtilecek \\
Program: & Lisans \\
Ana Bilim Dali: & Bilgisayar Muhendisligi \\
Teslim Tarihi: & April 2026 \\
\end{tabular}

\vfill
\begin{minipage}{0.92\textwidth}
\small
Not: Kose parantezli alanlar universite, ogrenci ve danisman bilgilerine gore doldurulmak uzere birakilmistir.
Bu surum akademik proposal biciminde, atifli ve kaynakcali olarak hazirlanmistir.
\end{minipage}
\end{center}
\end{titlepage}

\pagenumbering{roman}
\tableofcontents
\newpage

\section*{Ozet}
\addcontentsline{toc}{section}{Ozet}

Bu arastirma onerisi, klasik CPU scheduling literaturundeki \textit{adaptive time quantum secimi}
ile hardware reliability literaturundeki \textit{aging-aware runtime control} yaklasimlarini tek bir
kontrol cercevesinde birlestirmeyi amaclamaktadir. Mevcut Round Robin turevi scheduler calismalari,
quantum degerini burst time, ready-queue uzunlugu veya istatistiksel ozetler uzerinden belirlemeye
odaklanirken; reliability-aware sistemler termal yonetim, DVFS, task remapping ve online monitor geri
bildirimi kullanarak donanim omrunu korumaya calismaktadir \cite{behera2010,marr2022,srinivasan2004,kumar2011}.
Ancak scheduler seviyesinde quantum seciminin dogrudan thermal ve aging telemetry ile birlestirildigi
arastirma sayisi halen sinirlidir.

Bu calismada, queue-state, thermal durum ve aging estimate bilgisini ayni kontrol dongusunde kullanan
bir \textit{aging-aware adaptive scheduler} tasarlanmasi onerilmektedir. Onerilen yontemsel cerceve uc
asamalidir: (i) heuristic aging-penalty tabanli quantum secimi, (ii) physics-informed supervised learning
ile optimum quantum prediction, (iii) contextual bandit veya reinforcement learning tabanli policy ogrenimi.
Arastirmanin temel hipotezi, boyle bir scheduler'in ya ayni performans seviyesinde daha dusuk aging birikimi
ya da ayni aging butcesi altinda daha iyi performans saglayabilecegidir.

\textbf{Anahtar Kelimeler:} adaptive scheduling, Round Robin, time quantum, hardware aging,\\ reliability-aware control, reinforcement learning

\newpage
\pagenumbering{arabic}
\onehalfspacing

\section{Giris}

Time quantum secimi, zaman paylasimli scheduler tasariminin temel kararlarindan biridir. Quantum cok kucuk
secildiginde context switch maliyeti ve scheduler overhead artar; quantum cok buyuk secildiginde ise interaktiflik,
response time ve kisa islerin tamamlama davranisi olumsuz etkilenir. Bu nedenle dynamic Round Robin literaturu,
quantum degerinin workload'a gore uyarlanmasi gerektigini uzun suredir gostermektedir \cite{behera2010,noon2011,marr2022,zohora2024}.

Bununla birlikte scheduler kararlari yalnizca performans degil, ayni zamanda sistemin fiziksel sagligi uzerinde de
dolayli etkiler uretir. Dengesiz cekirdek kullanimi, hotspot olusumu, uzun sureli yuksek utilization ve termal
gradyanlar; NBTI, PBTI, HCI, electromigration ve thermal cycling gibi aging mekanizmalarini hizlandirabilir
\cite{srinivasan2004,agarwal2007,kumar2011}. Buna ragmen scheduling literaturu ile aging-aware control literaturu
buyuk olcude ayri gelismistir.

Bu proposal'in motivasyonu, bu iki damari birlestirmektir: scheduler, queue-state ile birlikte thermal ve aging
geri bildirimine de bakarak quantum, migration ve gerekiyorsa DVFS aksiyonlari uretebilir mi?

\section{Problem Tanimi}

Arastirmanin temel problemi su sekilde ifade edilebilir:

\begin{quote}
Sabit veya yalnizca queue-istatistiklerine dayali dynamic quantum secimi, scheduler kararlarinin uzun vadeli
hardware aging ve reliability maliyetlerini dikkate almamaktadir. Bu nedenle performans, fairness, termal durum
ve projected lifetime'i birlikte optimize eden yeni bir scheduler-level kontrol mekanizmasina ihtiyac vardir.
\end{quote}

Bu problem ozellikle multicore sistemlerde daha belirgindir; cunku farkli cekirdeklere gorev dagilimi, context-switch
ritmi ve termal yogunluk, aging davranisina zamansal ve mekansal olarak etki etmektedir.

\section{Arastirma Sorulari ve Hipotezler}

\subsection{Arastirma Sorulari}

\begin{enumerate}[leftmargin=*]
\item Queue-state ile thermal/aging telemetry birlikte kullanildiginda quantum secimi daha dengeli bir performance-reliability trade-off'u saglayabilir mi?
\item Aging-aware adaptive scheduler, sabit RR ve klasik dynamic RR yontemlerine gore ayni throughput seviyesinde daha dusuk aging score uretebilir mi?
\item Physics-informed supervised learning veya RL tabanli policy'ler, heuristic yontemlere gore workload degisimi altinda daha iyi genelleme yapabilir mi?
\item Scheduler karar uzayina migration ve DVFS eklendiginde projected lifetime uzerinde anlamli iyilesme elde edilir mi?
\end{enumerate}

\subsection{Hipotezler}

\begin{enumerate}[leftmargin=*]
\item Queue-aware ve aging-aware hibrit scheduler'lar, geleneksel dynamic RR scheduler'lara gore daha iyi bir latency-reliability dengesi saglar.
\item Thermal ve aging state quantum secimine dahil edildiginde hotspot davranisi ve birikimli stres dagilimi azalir.
\item Physics-informed supervised modeller, saf istatistik tabanli scheduler'lardan daha iyi quantum secimi yapabilir.
\item RL tabanli denetim, uygun reward ve guvenlik kisitlari ile uzun vadeli aging maliyetini azaltabilir.
\end{enumerate}

\section{Arastirmanin Onemi ve Katkisi}

Bu arastirmanin onemi dort noktada toplanabilir:

\begin{enumerate}[leftmargin=*]
\item Scheduling literaturu ile aging-aware reliability literaturunu birlestirmesi
\item Quantum secimini sadece performans degil, ayni zamanda donanim sagligi problemi olarak ele almasi
\item Heuristic, supervised learning ve RL tabanli yontemleri ortak bir simulator altinda karsilastirmasi
\item ``same performance, less aging'' ve ``same aging budget, better performance'' iddialarini nicel olarak test etmesi
\end{enumerate}

\section{Literatur Taramasi}

\subsection{Adaptive Time Quantum Scheduling}

Round Robin scheduler uzerinde quantum seciminin dinamiklestirilmesi uzun suredir calisilan bir alandir.
Behera ve ark. yeniden ayarlanmis dynamic quantum yontemi ile context switch ve turnaround time acisindan iyilesme
raporlamistir \cite{behera2010}. Noon ve ark. ise quantum degerini mean burst time tabanli bir kural ile secmistir
\cite{noon2011}. Daha yakin tarihli calismalarda median-average ve percentile tabanli daha dayanikli yaklasimlar
onerilmistir \cite{marr2022,zohora2024}. Bu literatur, workload-aware quantum seciminin faydali oldugunu gosterse de,
termal ve aging etkileri cogu zaman modelin disinda birakilmistir.

\subsection{AI Tabanli Scheduler Yaklasimlari}

AI tabanli scheduler calismalari artmakla birlikte, quantum secimini dogrudan bir sequential decision problemi olarak
ele alan calisma sayisi azdir. Chitaliya ve ark. quantum optimizasyonunu bir tahmin problemi olarak formulize edip
ML modelleri ile optimum quantum kestirimi denemistir \cite{chitaliya2023}. Buna karsilik scheduler-level RL calismalari
genis ama daginik bir literatur olusturmakta; gercek CPU quantum secimi ile dogrudan bag kuran calisma sayisi sinirlidir.

\subsection{Hardware Aging ve Reliability-Aware Control}

Srinivasan ve ark. lifetime reliability-aware microprocessors fikrini ortaya koyarak mikroislemci tasariminda
aging'i temel bir sistem seviyesi sorun olarak tanimlamistir \cite{srinivasan2004}. Agarwal ve ark. transistor aging
icin failure prediction mekanizmalarini one cikarirken \cite{agarwal2007}, Kumar ve ark. aging kaynakli performans
bozulumunu telafi etmek icin adaptif devre seviyesinde teknikler gelistirmistir \cite{kumar2011}. Reliability-aware
DVFS ve runtime thermal management calismalari da performans ile lifetime arasindaki trade-off'un scheduler ve runtime
katmaninda yonetilebilecegini gostermektedir \cite{aydin2009,yeganeh2021}.

\subsection{Literatur Boslugu}

Literaturdeki ana bosluklar su sekildedir:

\begin{itemize}[leftmargin=*]
\item Quantum secimi calismalari hardware aging bilgisini kullanmaz.
\item Aging-aware control calismalari scheduler quantum'unu dogrudan karar degiskeni olarak kullanmaz.
\item AI tabanli scheduler yaklasimlari reliability-centric objective kurmakta zayiftir.
\end{itemize}

Bu proposal tam olarak bu bosta konumlanmaktadir.

\section{Onerilen Yontem}

\subsection{Genel Cerceve}

Onerilen sistem dort temel bilesenden olusur:

\begin{enumerate}[leftmargin=*]
\item Workload generator veya trace source
\item Scheduler engine
\item Thermal estimator ve aging estimator
\item Logging, evaluation ve analysis katmani
\end{enumerate}

Policy katmani queue-state, thermal telemetry ve aging estimate bilgisini alarak quantum, migration ve opsiyonel
olarak DVFS karari uretecektir.

\subsection{Yontem 1: Heuristic Aging-Penalized Quantum}

Ilk yontem aciklanabilir bir baseline olarak tanimlanmaktadir:

\[
Q_t = f(S^{queue}_t) - \lambda_1 g(S^{thermal}_t) - \lambda_2 h(S^{aging}_t)
\]

Burada $f$, mean/median/quartile tabanli temel quantum uretir; $g$ ve $h$ ise sicaklik, hotspot severity,
thermal gradient, cumulative stress ve tahmini slack loss gibi sinyallere dayali ceza fonksiyonlaridir.

\subsection{Yontem 2: Physics-Informed Supervised Learning}

Bu asamada optimum quantum secimi bir prediction problemi olarak ele alinacaktir. Farkli workload pencereleri icin
farkli quantum degerleri taranarak bir objective tanimlanacaktir:

\[
J = w_1 \cdot \text{Latency} + w_2 \cdot \text{ContextSwitchCost} + w_3 \cdot \text{ThermalCost} + w_4 \cdot \text{AgingCost}
\]

Minimum $J$ veren quantum ``etiket'' olarak atanacak ve Random Forest, XGBoost, LightGBM gibi modellerle quantum
prediction yapilacaktir. Bu yontem physics-informed olarak adlandirilir; cunku etiketleme ve feature tasarimi thermal
ve aging anlami tasiyan degiskenler uzerinden kurulacaktir.

\subsection{Yontem 3: Contextual Bandit ve Reinforcement Learning}

Daha ileri asamada scheduler sequential decision problemi olarak ele alinacaktir. State uzayi queue length,
remaining burst statistics, per-core utilization, current temperature, temperature gradient, cumulative aging score,
onceki quantum ve migration history gibi ozellikleri icerecektir. Reward genel olarak asagidaki bicimde tasarlanabilir:

\[
R_t = \alpha \cdot \text{Throughput}_t - \beta \cdot \text{ResponseTime}_t - \gamma \cdot \text{ContextSwitches}_t - \delta \cdot \text{HotspotPenalty}_t - \eta \cdot \text{AgingIncrement}_t
\]

Bu asamada DQN, Double DQN veya PPO gibi algoritmalar denenebilir. Ancak tum AI policy'leri bir safety filter altinda
calisacak; peak temperature limiti, starvation prevention ve max migration rate gibi kurallar korunacaktir.

\section{Deneysel Tasarim}

\subsection{Simulator Ortami}

Ilk asamada Python tabanli discrete-event scheduler simulator kurulmasi planlanmaktadir. Bu secim, hizli deneysellik,
tekrarlanabilirlik ve policy varyantlarini kolay karsilastirma avantajlari saglamaktadir. Ileri asamada gem5 veya Sniper
entegrasyonu opsiyonel olarak dusunulebilir.

\subsection{Workload Siniflari}

\begin{longtable}{@{}p{3.3cm}p{2.1cm}p{8.0cm}@{}}
\caption{Planlanan workload siniflari.}\\
\toprule
\textbf{Sinif} & \textbf{Tur} & \textbf{Aciklama} \\
\midrule
\endfirsthead
\multicolumn{3}{l}{\small\itshape (devami)}\\
\toprule
\textbf{Sinif} & \textbf{Tur} & \textbf{Aciklama} \\
\midrule
\endhead
\bottomrule
\endfoot

Kisa CPU-bound & Sentetik & Kisa burst'ler; response time hassasiyeti yuksek \\
Uzun CPU-bound & Sentetik & Uzun burst'ler; quantum buyuklugune hassas \\
Heavy-tail bursts & Sentetik & Pareto benzeri uzun-kuyruk dagilim \\
Bursty arrivals & Sentetik & Kume halinde gelen prosesler ile queue baskisi \\
Periodic tasks & Sentetik & Gercek-zaman benzeri duzenli isler \\
Hot-core stress & Senaryo & Ayni cekirdekte aging birikimini zorlayan dagilim \\
Balanced multicore mix & Senaryo & Wear-leveling davranisini olcmeye uygun karisik yuk \\
Trace-driven mix & Gercekci & Mumkunse gercek scheduler trace'lerinden turetilmis isler \\
\end{longtable}

\subsection{Baseline Yontemler}

\begin{enumerate}[leftmargin=*]
\item Fixed Round Robin
\item Mean-based dynamic RR
\item Median/quartile-based dynamic RR
\item Thermal-aware scheduler
\item Aging-aware fakat AI icermeyen scheduler
\item Onerilen heuristic hybrid scheduler
\item Onerilen supervised scheduler
\item Onerilen bandit / RL scheduler
\end{enumerate}

\subsection{Thermal ve Aging Modelleri}

Termal model asagidaki sekilde basit ama tutarli bir formda ele alinacaktir:

\[
T_{t+1} = aT_t + bU_t + cS_t + dT_{amb}
\]

Burada $U_t$ utilization, $S_t$ switching proxy ve $T_{amb}$ ortam sicakligidir. Aging increment ise:

\[
\Delta A_t = k_1 e^{\rho T_t} + k_2 U_t^p + k_3 C_t
\]

Burada $C_t$, context-switch veya switching-rate kaynakli ek stresi temsil eder. Ileri asamada BTI-like,
HCI-like, EM proxy ve thermal-cycling bilesenleri ayri takip edilebilir.

\section{Degerlendirme Metrikleri}

\subsection{Scheduler Performansi}

\begin{itemize}[leftmargin=*]
\item Average waiting time
\item Average turnaround time
\item Average response time
\item Throughput
\item Context switch count
\item Fairness ve starvation gostergeleri
\end{itemize}

\subsection{Reliability ve Thermal Metrikleri}

\begin{itemize}[leftmargin=*]
\item Peak temperature
\item Mean temperature per core
\item Thermal gradient
\item Thermal cycling severity
\item Cumulative aging score
\item Projected lifetime proxy
\end{itemize}

\subsection{Temel Karsilastirma Iddialari}

Bu proposal su iki iddiayi test edecektir:

\begin{itemize}[leftmargin=*]
\item Ayni performans seviyesinde daha dusuk aging birikimi
\item Ayni aging butcesi altinda daha iyi performans
\end{itemize}

\section{Ornek Senaryo ve Objektif Beklenti Cercevesi}

Bu bolum, onerilen sistemin nasil davranacagini somut bir senaryo altinda gostermek ve ayni zamanda
calismadan beklenen fayda ile olasi sinirlari acik bicimde ifade etmek icin eklenmistir.

\subsection{Ornek Senaryo}

Dort cekirdekli bir sistem dusunulsun. Ready queue icinde hem kisa hem uzun CPU-bound isler bulunmaktadir.
Sistem belirli bir zaman penceresinde su kosullara girsin:

\begin{itemize}[leftmargin=*]
\item Core 1 uzun suredir yuksek utilization altinda calismaktadir.
\item Core 1 sicakligi diger cekirdeklerden daha yuksektir.
\item Core 1 icin aging score diger cekirdeklere gore daha hizli artmaktadir.
\item Queue doluluk orani yuksektir; dolayisiyla scheduler performans baskisi altindadir.
\end{itemize}

Bu durumda geleneksel bir dynamic RR scheduler buyuk olasilikla queue basincini azaltmak icin quantum'u arttiracak,
fakat bunu yaparken donanim sagligi bilgisini kullanmadigi icin ayni cekirdegin uzerindeki stresi daha da arttirabilecektir.
Onerilen aging-aware scheduler ise ayni anda su aksiyonlardan birini veya birkacini birlikte degerlendirecektir:

\begin{enumerate}[leftmargin=*]
\item quantum degerini queue baskisina gore arttirmak, fakat thermal/aging cezasi nedeniyle bu artis miktarini sinirlamak
\item isi daha az yorulmus bir cekirdege migrate etmek
\item gerekiyorsa ilgili cekirdekte DVFS ile frekans/voltaj seviyesini daha korumaci moda almak
\item fairness bozuluyorsa migration veya quantum kararini filtrelemek
\end{enumerate}

Bu senaryo, sistemin sadece ``hangi is ne kadar sure calissin'' degil, ayni zamanda ``bu karar donanimi nasil etkileyecek''
sorusuna cevap veren bir scheduler oldugunu gosterir.

\subsection{Muhtemel Pozitif Sonuclar}

Bu senaryo altinda beklenebilecek olumlu sonuclar sunlardir:

\begin{itemize}[leftmargin=*]
\item belirli cekirdeklerde hotspot davranisinin azalmasi
\item core bazli aging score dagiliminin daha dengeli hale gelmesi
\item sabit RR veya yalnizca queue-statistics tabanli RR'a gore benzer throughput altinda daha dusuk aging birikimi
\item thermal gradient ve peak temperature degerlerinde azalma
\item uzun sureli kosullarda projected lifetime proxy'sinde iyilesme
\end{itemize}

\subsection{Muhtemel Negatif Sonuclar ve Sinirlar}

Bu proposal'in objektif degerlendirilmesi icin, olumsuz veya sinirli sonuc ihtimalleri de acik bicimde kabul edilmelidir:

\begin{itemize}[leftmargin=*]
\item migration ve policy hesaplama overhead'i nedeniyle toplam latency bazen artabilir
\item aging korumasi fazla agresif olursa throughput veya turnaround time kotulesebilir
\item thermal bilgi aging'i tek basina temsil etmeyebilir; bu nedenle proxy modeli eksik kalabilir
\item RL tabanli policy, egitim verisine veya reward agirliklarina hassas olabilir
\item simulator sonucunda gozlenen fayda, gercek sistem entegrasyonunda daha kucuk kalabilir
\item bazi workload ailelerinde klasik dynamic RR zaten yeterince iyi calisiyorsa ek kazanc sinirli olabilir
\end{itemize}

\subsection{Objektif Basari Kriterleri}

Calisma ancak asagidaki turden bir sonuc gosterebilirse guclu kabul edilmelidir:

\begin{enumerate}[leftmargin=*]
\item ayni throughput altinda istatistiksel olarak anlamli aging score azalmasi
\item ayni aging butcesi altinda waiting time veya turnaround time iyilesmesi
\item peak temperature ve hotspot sayisinda acik azalma
\item birden fazla workload sinifinda tutarli iyilesme
\end{enumerate}

Eger bu kriterler saglanmazsa, proposal'in degerli ciktisi yine de su olabilir: scheduler-level aging-aware kontrolun
hangi kosullarda faydali oldugu ve hangi kosullarda overhead nedeniyle zayif kaldiginin acikca ortaya konmasi.

\section{Beklenen Zorluklar ve Risk Azaltma}

\begin{table}[h!]
\centering
\renewcommand{\arraystretch}{1.3}
\begin{tabularx}{\textwidth}{lXX}
\toprule
\textbf{Risk} & \textbf{Sorun} & \textbf{Azaltma Stratejisi} \\
\midrule
Aging model sade kalmasi & Fiziksel gercekligi eksik temsil etme & Sensitivity analysis, alternatif proxy'ler ve parametre taramasi \\
Sentetik workload bias & Gercek sistem davranisinin eksik yansimasi & Trace-driven workload ve farkli dagilim aileleri kullanma \\
RL kararsizligi & Reward hassasiyeti ve egitim zorlugu & Once heuristic ve supervised baseline'lari oturtma \\
Policy overhead & Kontrol maliyetinin kazanci silmesi & Ayrik action space ve seyrek update \\
Fairness bozulmasi & Aging korumasi bazi isleri cezalandirabilir & Fairness cezasi ve starvation bounds \\
\bottomrule
\end{tabularx}
\end{table}

\section{Calisma Plani}

\begin{center}
\begin{ganttchart}[
    hgrid,
    vgrid,
    time slot format=isodate-yearmonth,
    x unit=0.95cm,
    y unit title=0.6cm,
    y unit chart=0.7cm,
    title/.append style={fill=deepblue!10},
    bar/.append style={fill=accentblue!60},
    bar height=0.55
]{2026-05}{2026-11}
\gantttitlecalendar{year, month}\\
\ganttbar{Literatur ve proposal sonlandirma}{2026-05}{2026-05}\\
\ganttbar{Baseline simulator}{2026-05}{2026-06}\\
\ganttbar{Thermal ve aging model entegrasyonu}{2026-06}{2026-07}\\
\ganttbar{Heuristic hybrid scheduler}{2026-07}{2026-08}\\
\ganttbar{Veri uretimi ve etiketleme}{2026-08}{2026-09}\\
\ganttbar{Supervised model deneyleri}{2026-09}{2026-10}\\
\ganttbar{Bandit / RL deneyleri}{2026-10}{2026-11}\\
\ganttbar{Ablation, analiz ve final yazim}{2026-11}{2026-11}
\end{ganttchart}
\end{center}

\section{Beklenen Sonuclar}

Beklenen sonuc, onerilen scheduler'in geleneksel scheduler'lara gore daha iyi bir performance-reliability
trade-off'u gostermesidir. En guclu deneysel sonuc tipleri sunlar olacaktir:

\begin{enumerate}[leftmargin=*]
\item Benzer throughput altinda anlamli aging score azalmasi
\item Benzer aging limiti altinda daha dusuk waiting/turnaround time
\item Peak temperature ve hotspot davranisinda iyilesme
\item Workload aileleri boyunca daha tutarli scheduler davranisi
\end{enumerate}

Ancak objektif bir degerlendirme icin su ihtimaller de raporlanmalidir:

\begin{itemize}[leftmargin=*]
\item bazi workload'larda kazanc cok kucuk olabilir
\item bazi policy'ler aging'i azaltirken latency'yi arttirabilir
\item en iyi performansi veren policy ile en iyi reliability sonucu veren policy ayni olmayabilir
\item RL tabanli surum, heuristic veya supervised surumden her zaman daha iyi cikmayabilir
\end{itemize}

Dolayisiyla calismanin amaci yalnizca ``onerilen yontem her kosulda en iyidir'' demek degil; hangi kosullarda avantaj
sagladigini, hangi kosullarda maliyet olusturdugunu ve trade-off yuzeyinin nasil sekillendigini gostermektir.

\section{Sonuc}

Bu proposal, adaptive time quantum scheduling ile aging-aware runtime control arasindaki boslugu doldurmaya yonelik
ayrintili ve uygulanabilir bir arastirma cercevesi sunmaktadir. Onerilen sistem, quantum secimini yalnizca performans
problemine indirgememekte; ayni zamanda donanim sagligini da karar mekanizmasina dahil etmektedir. Bu yonuyle calisma,
isletim sistemleri, bilgisayar mimarisi ve AI for systems alanlarinin kesisiminde anlamli bir katki potansiyeli tasimaktadir.

\newpage
\begin{thebibliography}{99}

\bibitem{behera2010}
H. S. Behera, R. K. Mohanty, and D. Nayak,
``A New Proposed Dynamic Quantum with Re-Adjusted Round Robin Scheduling Algorithm and Its Performance Analysis,''
\textit{International Journal of Computer Applications}, 2010.

\bibitem{noon2011}
A. Noon, A. Kalakech, and S. Kadry,
``A New Round Robin Based Scheduling Algorithm for Operating Systems: Dynamic Quantum Using the Mean Average,''
arXiv preprint arXiv:1111.5348, 2011.

\bibitem{marr2022}
Sakshi, C. Sharma, S. Sharma, S. Kautish, S. A. M. Alsallami, E. M. Khalil, and A. W. Mohamed,
``A New Median-Average Round Robin Scheduling Algorithm,''
\textit{Alexandria Engineering Journal}, 2022.

\bibitem{zohora2024}
M. F. Zohora, F. Farhin, and M. S. Kaiser,
``An Enhanced Round Robin Using Dynamic Time Quantum for Real-Time Asymmetric Burst Length Processes,''
\textit{PLOS ONE}, 2024.

\bibitem{chitaliya2023}
P. Chitaliya et al.,
``Time Quantum Optimization in Round Robin Algorithm,''
\textit{NMITCON}, 2023.

\bibitem{srinivasan2004}
J. Srinivasan et al.,
``The Case for Lifetime Reliability-Aware Microprocessors,''
\textit{ISCA}, 2004.

\bibitem{agarwal2007}
M. Agarwal et al.,
``Circuit Failure Prediction and Its Application to Transistor Aging,''
\textit{VTS}, 2007.

\bibitem{aydin2009}
H. Aydin and D. Zhu,
``Reliability-Aware Energy Management for Periodic Real-Time Tasks,''
\textit{IEEE Transactions on Computers}, 2009.

\bibitem{kumar2011}
S. V. Kumar, C. H. Kim, and S. S. Sapatnekar,
``Adaptive Techniques for Overcoming Performance Degradation Due to Aging in CMOS Circuits,''
\textit{IEEE Transactions on VLSI Systems}, 2011.

\bibitem{yeganeh2021}
A. Yeganeh-Khaksar et al.,
``Ring-DVFS: Reliability-Aware Reinforcement Learning-Based DVFS for Real-Time Embedded Systems,''
\textit{IEEE Embedded Systems Letters}, 2021.

\end{thebibliography}

\end{document}
